\section{Conclusion}\label{sec:conclusion}
This study examined local texture evidence and spatial aggregation in a compact tyre-inspection system using existing images and preserved training records. Dense annotation agreement coexisted with poor mileage-proxy transfer across the allocated tyres. TRACE subsequently improved fixed local and masked descriptor reconstruction, but no epoch met its original global-diversity requirements.\AIdraftingnote

The saved-feature analysis showed that a global mean retained 61.14\% of local validation advantage over the position control, compared with 81.76\% for a position template plus five spatial residual moments. Descriptor losses and component summaries showed that this retention was not uniform. A fixed-mass covariance partition, invertible scaling control, and exploratory mask-layout CKA further separated useful local readout behavior from spectral diversity.

The supported conclusion is that local evidence can remain useful while the global representation is inadequate under its declared contract. Preserving local spatial residuals is a justified next research question, not a proven solution to proxy transfer. The current endpoint remains unselected, and the evidence does not establish physical wear, depth, alignment, or safety accuracy. The complete failure record and explicit computation boundary make this limited-data result inspectable without another training run.
